\chapter{Termination and Reduction Orderings}
\label{ch:rewrite-termination}

An optimizer that repeatedly applies individually valid rules can run forever.
Commuting two operands back and forth preserves meaning at every step yet never finishes.
Termination therefore needs a decreasing global argument, not only semantic correctness of each local replacement.
The argument must cover substitutions and contexts, because the rule engine applies schematic rules inside larger programs.

\section{Find a measure that every step decreases}
\label{sec:ch83-measure}
For strings, rules with strictly shorter right sides decrease the natural-number length measure.
For finite trees, removing add-zero and multiply-one nodes decreases tree size for every substitution.
Let $|x|=1$ for a variable and $|f(t_1,\ldots,t_n)|=1+\sum_i|t_i|$.
The rule $\operatorname{add}(x,0)\to x$ removes two nodes regardless of the size substituted for x.

But a symbolic size comparison on schematic terms alone is insufficient when variable multiplicities change.
The rule $f(x)\to g(x,x)$ duplicates the substituted term.
Assigning the symbol f a large constant weight can make the schematic left look heavier for x of size one while failing for a sufficiently large substitution.
Table~\ref{tab:ch83-measures} lists the obligations that a usable measure must satisfy.

\begin{table}[htbp]
\caption{A local decrease must survive every substitution and surrounding context.}
\label{tab:ch83-measures}
\centering
\begin{tabular}{ll}
\toprule
Property & Required role \\
\midrule
well-founded strict order & excludes infinite descending chains \\
substitution stability & schematic decrease holds for instantiated rules \\
context monotonicity & replacing a smaller subterm decreases the whole \\
strict orientation of all rules & every legal rewrite step is covered \\
\bottomrule
\end{tabular}
\end{table}

A \term{reduction ordering} is a well-founded strict order stable under substitution and contexts.
If every rule decreases it, every step decreases it, so an infinite rewrite sequence would imply an impossible infinite descent~\cite{dershowitz-1982-orderings}.
This proves termination of the specified relation.
It does not prove confluence, an efficient bound on derivation length, or semantic validity.

\section{Use polynomial interpretations with explicit domains}
\label{sec:ch83-polynomial}
An interpretation maps each symbol to a monotone function over a well-founded domain, commonly natural numbers.
Variables range over arbitrary domain values.
For the identity rules, interpret add and mul structurally as $1+a+b$ and zero and one as $1$; these are measures, not the arithmetic meanings of the symbols.
Then add(x,zero) has measure $x+2$, strictly greater than x for every natural x.
The same holds for mul(x,one).

Figure~\ref{fig:ch83-potential} shows the potential decrease for the nested worked term.
If v has measure one, the initial term has measure five, the intermediate identity wrapper has measure three, and the result has measure one.
Every identity step drops two, so this particular family has exactly one step per wrapper.
The tests construct up to 49 pairs of wrappers and check both trace length and normal form.

\begin{figure}[htbp]
\centering
\begin{tikzpicture}[node distance=14mm]
\node[basic block] (a) {mul(add(v,zero),one): 5};
\node[basic block,below=of a] (b) {add(v,zero): 3};
\node[basic block,below=of b] (c) {v: 1};
\draw[cfg edge] (a)--node[right,font=\small]{decrease 2}(b);
\draw[cfg edge] (b)--node[right,font=\small]{decrease 2}(c);
\end{tikzpicture}
\caption{A concrete potential can bound the trace when every step has a known minimum decrease.}
\label{fig:ch83-potential}
\end{figure}

Algorithm~\ref{alg:ch83-linear} checks a restricted affine interpretation.
For a difference $c+\sum_i a_i x_i$ over natural variables, nonnegative coefficients and $c\ge1$ suffice for strict positivity.
Requiring every symbol's argument coefficients to be positive supplies strict context monotonicity.
These are sufficient syntactic checks for this domain, not a complete termination prover.

\begin{algorithm}[htbp]
\caption{Check a sufficient affine interpretation certificate over natural numbers.}
\label{alg:ch83-linear}
\begin{pseudocode}
\Require finite rules and nonnegative affine symbol interpretations
\Ensure accepted certificate proves every rule strictly decreasing, or unresolved
\State require each argument coefficient to be positive
\ForAll{rules $l\to r$}
\State compose interpretations to obtain affine forms for $l$ and $r$
\State subtract right coefficients from left coefficients
\State require every variable coefficient nonnegative and constant difference positive
\EndFor
\State accept under the stated natural-number interpretation
\end{pseudocode}
\end{algorithm}

With $V$ variables and $N$ total term nodes, direct coefficient-vector composition costs $O(NV)$ arithmetic operations; integer coefficient bit sizes add further cost.
Failure means the proposed interpretation does not establish the decrease through this sufficient check.
It does not establish nontermination.
Searching coefficients requires a separate finite search or solver whose completeness and resource behavior must be stated.

\section{Orient rules that preserve or increase size}
\label{sec:ch83-path}
Many useful rules do not shrink plain node count.
Changing a representation can move a complex operator under simpler operators, or flatten a nested form without removing symbols.
\term{Recursive path orderings} compare root precedence and recursively compare arguments while respecting the subterm property~\cite{dershowitz-1982-orderings}.
Their power comes from structural descent, not a claim that every source-level algebraic identity can be oriented compatibly.

A precedence $f\succ g$ can justify $f(x)\succ g(x)$ under an appropriate path ordering, while the subterm property gives $f(x)\succ x$.
Comparing $f(a)$ with $g(b)$ also requires the relevant dominance checks on arguments; root precedence alone is not enough.
Lexicographic versus multiset statuses change how argument tuples are compared.
A production ordering implementation needs precise rules for variables, equal roots, statuses, and precedence ties before accepting a certificate.

The chapter does not implement a general recursive path ordering checker.
It uses the ordering literature to explain alternatives when the executed size-decreasing examples are insufficient.
Similarly, dependency-pair methods separate potential recursive calls from surrounding construction and can prove systems beyond simple order orientation; their usable-rule, reduction-pair, and graph approximations require their own correctness premises~\cite{ohlebusch-2002-rewriting}.
They should not be reduced to a slogan that finding no obvious syntactic cycle proves termination.

\section{Distinguish semantic recursion from engine looping}
\label{sec:ch83-engine}
Consider commutativity $\operatorname{add}(x,y)\to\operatorname{add}(y,x)$.
For unequal operands, the rule can immediately undo itself.
Sorting operands by a stable total order would instead orient only out-of-order pairs, decreasing an inversion-like measure in a fixed associative representation.
That implementation also needs to state whether reassociation and operand reordering preserve effects and arithmetic semantics.

Table~\ref{tab:ch83-loop} contrasts ways to control repeated rewriting.
A fuel limit guarantees the engine returns control, not that the mathematical system terminates.
A visited-state set can detect a repeated state on a deterministic path, proving a cycle in that path, but unbounded growth can diverge without repeating any state.
A node-size limit bounds memory only if it accounts for all retained structures, not just the current root.

\begin{table}[htbp]
\caption{Resource controls and termination proofs answer different questions.}
\label{tab:ch83-loop}
\centering
\begin{tabular}{ll}
\toprule
Mechanism & What it establishes \\
\midrule
well-founded decrease & no infinite reduction under covered rules \\
fuel & at most the configured number of performed steps \\
repeated-state detection & a cycle when a state actually repeats \\
term-size cap & bounded admitted term size, not normal-form existence \\
heuristic rule phases & selected schedule, requiring separate termination analysis \\
\bottomrule
\end{tabular}
\end{table}

Algorithm~\ref{alg:ch83-bounded} is a safe publication wrapper for a budgeted normalizer.
It must test redex absence after the final permitted step: using exactly all fuel does not imply incompleteness if the last step reached a normal form.
The companion implements that final test and returns its status separately from the current term.

\begin{algorithm}[htbp]
\caption{Report normalization status independently of the number of performed steps.}
\label{alg:ch83-bounded}
\begin{pseudocode}
\Require finite term, valid rules, nonnegative step budget
\Ensure term, trace, and accurate normal-form status
\State perform selected steps while both a redex and budget remain
\State test whether the final term has any selected-strategy redex
\If{no redex remains} \State return complete normal form and trace
\Else \State return intermediate term, trace, and incomplete status
\EndIf
\end{pseudocode}
\end{algorithm}

The wrapper costs one additional redex search beyond the performed steps in the exhausted case.
For a strategy intended to normalize the full relation, the final search must cover all allowed positions and rules.
A phase that deliberately ignores some positions can only report normality relative to that phase.
Figure~\ref{fig:ch83-outcomes} keeps these conclusions distinct.

\begin{figure}[htbp]
\centering
\begin{tikzpicture}[node distance=13mm and 14mm]
\node[basic block] (run) {bounded rewrite run};
\node[basic block,below left=of run] (done) {no redex: normal form};
\node[basic block,below right=of run] (stop) {redex remains: incomplete};
\draw[cfg edge] (run)--(done);\draw[cfg edge] (run)--(stop);
\end{tikzpicture}
\caption{Budget exhaustion must not be relabeled as a proof of normalization.}
\label{fig:ch83-outcomes}
\end{figure}

\section{Construct the descent proof and its scope}
\label{sec:ch83-proof}
Assume a strict well-founded ordering greater-than that is stable under substitution and monotone under every admitted context.
If l is greater than r for each rule, then an instance l substituted by sigma is greater than r substituted by sigma.
Putting those two terms into the same context preserves the comparison.
Thus each actual engine step decreases the whole represented term, not merely the matched fragment.
An infinite reduction would induce an infinite descending chain, contradicting well-foundedness.
This short argument identifies exactly where the substitution and context premises are used.

It also explains why a heuristic cost model is usually not a termination certificate.
Estimated target costs can decrease locally while context-sensitive sharing or register pressure increases globally.
Costs can be negative, incomparable, or revised by a learned estimator between iterations.
A stable syntactic potential used solely to control rewrite progress can coexist with a separate profitability model, but they should not be confused.
The progress measure need not predict execution speed, and the speed estimate need not be well-founded.

For a natural-number measure with minimum step decrease one, derivation length is at most the initial measure minus the minimum reachable measure.
This yields a useful bound for the executed identity wrappers.
It does not yield a polynomial source-size bound for every termination ordering: a measure can be enormous relative to input encoding, and intermediate terms can grow while a lexicographic or structural ordering decreases.
The time per step can also grow with those terms.
Termination, derivational complexity, and practical optimizer latency remain three separate conclusions.

\section{Combine passes without introducing an oscillation}
\label{sec:ch83-composition}
Two terminating passes can form a nonterminating repeated pipeline.
One pass expands a representation and terminates because it visits each original node once; another contracts the expanded representation and also terminates.
Repeating both until neither changes anything can recreate the same input forever.
Each pass's internal proof says nothing about the composite loop unless the proofs share a global measure or the pipeline has an independent phase bound.

A compiler can deliberately permit one expansion phase followed by one contraction phase and stop.
It can also attach a bounded phase count, remember previously seen pipeline states, or require a lexicographic potential covering all phases.
Repeated-state detection catches actual cycles but not indefinitely growing oscillations whose states never repeat.
If the compiler's observable output must be deterministic, budget accounting and tie-breaking should not depend on unstable hash iteration order.

Recursive and conditional rules require equally careful scope.
A termination certificate for an unconditional overapproximation proves termination of any restriction of that relation, including rules whose conditions disable some steps.
But a certificate for one restricted strategy does not automatically prove termination when the engine starts rewriting under lambdas or inside previously opaque nodes.
The set of contexts allowed by the engine is part of the theorem.
Changing it is a semantic and termination-policy change, even if the rule table itself is byte-for-byte identical.

\section{Exercises}
\label{sec:ch83-exercises}
\begin{enumerate}
\item\label{ex:ch83-size} Compute the size decrease of add(x,zero)$\to$x after substituting a term of size 17.
\item\label{ex:ch83-duplicate} With interpretations $f(x)=10+x$ and $g(x,y)=1+x+y$, determine for which natural x the rule $f(x)\to g(x,x)$ decreases.
\item\label{ex:ch83-fuel} Normalize add(v,zero) with one step of fuel.
Should its final status be complete? Contrast a self-loop rule.
\item\label{ex:ch83-commute} Give the two-step cycle generated by unconditional add commutativity and name an additional canonical-order policy that blocks it.
\end{enumerate}

\section{Further reading}
\label{sec:ch83-reading}
Dershowitz develops simplification orderings and recursive path orderings~\cite{dershowitz-1982-orderings}.
Ohlebusch treats more advanced termination techniques, including the conditions that make modular proof methods sound~\cite{ohlebusch-2002-rewriting}.
